\section{Methodology}

Our approach uses recursive bisection to partition census tracts into congressional districts. We present two modes: \textbf{baseline recursive bisection} using unweighted edges, and \textbf{edge-weighted recursive bisection} using boundary lengths as edge weights. Both modes satisfy constitutional requirements (equal population, contiguity); edge-weighted mode additionally optimizes compactness.

\subsection{Problem Formulation}

Given a state with population $P$ allocated $k$ congressional districts:

\textbf{Inputs:}
\begin{itemize}
\item Census tracts $V = \{v_1, \ldots, v_n\}$ with populations $p_i$
\item Adjacency graph $G = (V, E)$ where $(i,j) \in E$ if tracts $i$ and $j$ share a boundary
\item Boundary lengths $w_{ij}$ for each edge $(i,j) \in E$ (edge-weighted mode only)
\end{itemize}

\textbf{Objective:}
Partition $V$ into $k$ disjoint districts $D_1, \ldots, D_k$ such that:

\begin{enumerate}
\item \textbf{Equal population}: Each district has population $\approx P/k$ within tolerance
\item \textbf{Contiguity}: Each district forms a connected subgraph
\item \textbf{Compactness} (edge-weighted mode): Minimize total perimeter $\sum_{d=1}^k \text{Perimeter}(D_d)$
\end{enumerate}

Target population per district: $\text{target}_d = P/k$. We allow $\pm 0.5\%$ deviation to balance districts precisely while giving METIS flexibility for compactness optimization.

\subsection{Baseline Recursive Bisection}

\subsubsection{Algorithm Overview}

Recursive bisection hierarchically divides regions into smaller subregions until reaching individual districts:

\begin{algorithm}[H]
\caption{Baseline Recursive Bisection}
\begin{algorithmic}[1]
\Function{Partition}{$G, k$}
    \If{$k = 1$}
        \State \Return single district containing all vertices
    \EndIf
    \State $k_1 \gets \lfloor k/2 \rfloor$
    \State $k_2 \gets k - k_1$
    \State Compute target populations $t_1, t_2$ for split
    \State $(V_1, V_2) \gets \Call{METIS-Bisect}{G, t_1, t_2}$
    \State $D_1 \gets \Call{Partition}{G[V_1], k_1}$
    \State $D_2 \gets \Call{Partition}{G[V_2], k_2}$
    \State \Return $D_1 \cup D_2$
\EndFunction
\end{algorithmic}
\end{algorithm}

At each bisection, METIS divides the current region into two parts with specified populations. This process continues until each leaf contains one district. The recursive structure creates a binary tree of splits, providing a clear decomposition hierarchy.

\subsubsection{METIS Bisection (Unweighted)}

We invoke METIS with:
\begin{itemize}
\item \texttt{-contig}: Enforce contiguity (connected partitions)
\item \texttt{-minconn}: Minimize connectivity (prefer cleaner boundaries)
\item \texttt{-ufactor=1.005}: Allow 0.5\% population imbalance
\item \texttt{-niter=100}: Refinement iterations for quality
\end{itemize}

METIS minimizes edge cuts: $\min \sum_{(i,j) \in E} \mathbb{1}[P_i \neq P_j]$ subject to population balance. This treats all adjacencies equally—cutting a short edge costs the same as cutting a long edge.

\subsubsection{Graph Representation}

Census tracts form graph vertices. Two tracts are adjacent if they share any boundary (Queen contiguity: edge or vertex contact). For states with islands or water barriers (30+ states including Hawaii, Alaska, Washington, Michigan, New York, California, Florida, and Maine), we use county-based bridge connections to ensure graph connectivity. For each disconnected component, we connect tracts to the nearest tract in the main component within the same county (using the first 5 digits of the 11-digit GEOID: SSCCCTTTTTT where SS = state FIPS, CCC = county FIPS). This minimal bridging approach creates only necessary water crossings while respecting county boundaries, unlike distance-threshold methods which create excessive edges or fail for distant islands.

Graph files use METIS CSR format:
\begin{verbatim}
<n_vertices> <n_edges> 001 1
<vweight_1> <neighbor_1> <neighbor_2> ...
<vweight_2> <neighbor_1> <neighbor_2> ...
...
\end{verbatim}

Format code 001 indicates vertex weights only (populations). Edges are unweighted, implicitly having weight 1.

\subsection{Edge-Weighted Recursive Bisection}

\subsubsection{Motivating Example}

Consider three adjacent tracts:
\begin{itemize}
\item Tract A and Tract B share a 150-meter boundary
\item Tract A and Tract C share a 12,500-meter boundary
\item Tract B and Tract C share a 200-meter boundary
\end{itemize}

Baseline recursive bisection treats all edges equally: cutting $(A,B)$, $(A,C)$, or $(B,C)$ has identical cost. METIS may separate A from C, adding 12.5 km to total perimeter.

Edge-weighted recursive bisection weights edges by actual length: $w_{AB} = 150$, $w_{AC} = 12500$, $w_{BC} = 200$. Now METIS strongly prefers cuts minimizing boundary length. Separating A from C costs 83× more than separating A from B. This directly optimizes geometric compactness.

\subsubsection{Boundary Length Computation}\label{sec:edge-weights}

For each adjacent tract pair $(i,j)$, we compute shared boundary length $w_{ij}$:

\begin{enumerate}
\item \textbf{Load geometries}: Read census tract shapefiles (Census TIGER/Line format)
\item \textbf{Project to CRS}: Transform to state-specific projected coordinate system (e.g., California Albers EPSG:3310) for accurate length measurements in meters
\item \textbf{Compute intersection}: $I_{ij} = \partial G_i \cap \partial G_j$ where $\partial G_i$ is tract $i$'s boundary
\item \textbf{Classify intersection}:
\begin{itemize}
    \item \textbf{LineString/MultiLineString}: $w_{ij} = \text{length}(I_{ij})$ (real shared boundary)
    \item \textbf{Point}: $w_{ij} = 0.1$ meters (corner contact)
    \item \textbf{Empty}: $w_{ij} = \text{median}\{w_{kl} \mid (k,l) \in E_{\text{land}}\}$ (county-based bridge)
\end{itemize}
\end{enumerate}

County-based bridge edges (water crossings created to connect disconnected components) use the median land boundary length for the state. This provides adaptive scaling: small states have shorter typical boundaries ($\sim$1 km), large states longer ($\sim$5-10 km). This balances permitting necessary water crossings while discouraging excessive splits. Critically, these bridges only connect tracts within the same county, ensuring geographic sensibility even for distant islands like Alaska's Aleutians ($\sim$1700 km separation) or Hawaii's inter-island gaps ($\sim$100 km).

\textbf{Computational complexity}: $O(Nd)$ for $N$ tracts with average degree $d \approx 6$. Preprocessing takes 10-30 seconds per state. Edge weights are cached as Python dictionaries, requiring $\sim$100 KB per 1,000 edges.

\subsubsection{METIS Integration with Edge Weights}

METIS supports edge weights via format code 011:
\begin{verbatim}
<n_vertices> <n_edges> 011 1
<vweight_1> <nbr_1> <eweight_1> <nbr_2> <eweight_2> ...
<vweight_2> <nbr_1> <eweight_1> <nbr_2> <eweight_2> ...
...
\end{verbatim}

Boundary lengths (in meters) are scaled to integer centimeters: $w_{ij}^{\text{METIS}} = \max(1, \lfloor w_{ij} \times 100 \rfloor)$. This provides 1 cm precision while avoiding numerical issues with small weights.

METIS now minimizes weighted edge cuts:
\[
\min_{P} \sum_{(i,j) \in E} w_{ij} \cdot \mathbb{1}[P_i \neq P_j]
\]

subject to population balance and contiguity. Since $w_{ij}$ represents boundary length, this directly minimizes total perimeter.

\subsubsection{Polsby-Popper Optimization}

The Polsby-Popper score for a district with area $A$ and perimeter $P$ is:
\[
PP = \frac{4\pi A}{P^2}
\]

For redistricting, area is determined by population (held constant by equal population constraint). Thus minimizing $P$ maximizes $PP$. Edge-weighted recursive bisection optimizes Polsby-Popper by construction.

\textbf{Global vs. local optimization}: Recursive bisection makes locally optimal splits at each level. This is not guaranteed globally optimal (NP-hard) but produces high-quality solutions in practice. The hierarchical structure also provides intuitive geographic divisions.

\subsection{Visualization of Recursive Bisection Process}

To illustrate the edge-weighted recursive bisection process, we present two examples showing how states are progressively divided.

\textbf{Minnesota (Even Division)}: Minnesota receives 8 congressional districts, an even power of 2. Figures~\ref{fig:minnesota_rounds} shows the three bisection rounds:
\begin{itemize}
\item \textbf{Round 1}: State divided into 2 regions (4 districts each)
\item \textbf{Round 2}: Each region bisected into 4 total regions (2 districts each)
\item \textbf{Round 3}: Final bisection creates 8 districts (1 district each)
\end{itemize}

Every split is perfectly balanced (equal-sized partitions), creating a symmetric partition tree. This even division allows maximum flexibility for METIS to optimize compactness at each level.

\begin{figure*}[t]
\centering
\begin{subfigure}[b]{0.32\textwidth}
    \includegraphics[width=\textwidth]{figures/minnesota_round_1_2_regions.png}
    \caption{Round 1: 2 regions}
\end{subfigure}
\hfill
\begin{subfigure}[b]{0.32\textwidth}
    \includegraphics[width=\textwidth]{figures/minnesota_round_2_4_regions.png}
    \caption{Round 2: 4 regions}
\end{subfigure}
\hfill
\begin{subfigure}[b]{0.32\textwidth}
    \includegraphics[width=\textwidth]{figures/minnesota_round_3_8_regions.png}
    \caption{Round 3: 8 districts}
\end{subfigure}
\caption{Minnesota edge-weighted recursive bisection progression. Three rounds of balanced bisection create 8 highly compact districts. Each split minimizes boundary length, producing districts with smooth, regular boundaries.}
\label{fig:minnesota_rounds}
\end{figure*}

\textbf{Alabama (Odd Division)}: Alabama receives 7 congressional districts, not a power of 2. Figure~\ref{fig:alabama_rounds} shows the three bisection rounds:
\begin{itemize}
\item \textbf{Round 1}: State divided into 2 regions (4 and 3 districts respectively)
\item \textbf{Round 2}: Unequal split creates 4 total regions (2, 2, 2, and 1 districts)
\item \textbf{Round 3}: Final bisection creates 7 districts (some regions split, some remain)
\end{itemize}

Odd division creates asymmetric partition trees—some regions are subdivided more than others. This is fundamental to recursive bisection with non-power-of-2 targets. The algorithm maintains population balance despite asymmetry by adjusting target populations at each split.

\begin{figure*}[t]
\centering
\begin{subfigure}[b]{0.32\textwidth}
    \includegraphics[width=\textwidth]{figures/alabama_round_1_2_regions.png}
    \caption{Round 1: 2 regions}
\end{subfigure}
\hfill
\begin{subfigure}[b]{0.32\textwidth}
    \includegraphics[width=\textwidth]{figures/alabama_round_2_4_regions.png}
    \caption{Round 2: 4 regions}
\end{subfigure}
\hfill
\begin{subfigure}[b]{0.32\textwidth}
    \includegraphics[width=\textwidth]{figures/alabama_round_3_7_regions.png}
    \caption{Round 3: 7 districts}
\end{subfigure}
\caption{Alabama edge-weighted recursive bisection progression. Asymmetric splits accommodate 7 districts (non-power-of-2). Despite unbalanced partition tree, edge-weighted mode produces compact districts through boundary length minimization.}
\label{fig:alabama_rounds}
\end{figure*}

These visualizations demonstrate that edge-weighted recursive bisection produces compact districts regardless of even/odd division. The key is boundary length minimization at each split—METIS prefers cuts that create smooth, regular boundaries rather than jagged, irregular ones.

\subsection{Algorithmic Properties}

\textbf{Determinism}: Given fixed METIS random seed, the algorithm produces identical results. This ensures reproducibility.

\textbf{Transparency}: The complete partition hierarchy is recorded. Every district's formation can be traced through the binary tree of splits, as shown in Figures~\ref{fig:minnesota_rounds} and~\ref{fig:alabama_rounds}.

\textbf{Computational efficiency}: Runtime is $O(Nk \log k)$ for $N$ tracts and $k$ districts. Typical states partition in 1-5 minutes.

\textbf{Constraints}: Both modes guarantee equal population (within 0.5\%) and contiguity. No additional constraints on communities of interest, political boundaries, or partisan balance are imposed.

\textbf{Fairness}: The algorithm treats all census tracts identically. No geographic regions receive preferential treatment. The only inputs are population counts and spatial adjacency—no political or racial data.
